% Generated by tools/textbook/export_listings.js from decks/plan.js and code/l19.S: do not edit.
\section{Il programma di questo capitolo}\label{sec:lab:l19}

I listati di questo capitolo sono parti di un solo programma,
\fname{code/l19.S} nel repository del corso, stampato per intero nel
Listato~\ref{lst:lab:l19}. I test automatici del materiale del corso lo
compilano e lo eseguono sull'esecutore senza interfaccia di \prgqt{}, con
ingressi programmati, e controllano il frame; la
Figura~\ref{fig:lab:l19} è il frame scritto dal test.

\begin{description}[style=nextline,leftmargin=1.5em]
\item[Sullo schermo.] Quattro righe di 32 lampade, il codice macchina di lw t1, 0(t0) / addi t1, t1, 1 / sw t1, 0(t0) / ret, e accanto a ciascuna un quadrato verde: la parola in memoria è quella decodificata a mano a lezione.
\item[Che cosa premere.] Niente: il programma va da solo.
\end{description}

\begin{figure}[htb]
\centering
\includegraphics[width=0.5\textwidth]{lab_l19.png}
\caption{\fname{l19.S} su \prgqt{}.}\label{fig:lab:l19}
\end{figure}

\asmfile[caption={Il programma completo del capitolo, \fname{code/l19.S}.},label=lst:lab:l19]{code/l19/program.S}

Per eseguirlo su \prgqt{} i comandi sono quelli del
Listato~\ref{lst:lab:l19:run}; il Paragrafo~\ref{sec:prg:testing} dà gli
stessi passi per QEMU e per la scheda.

\begin{shell}[caption={Compilare \fname{l19.S} ed eseguirlo su \prgqt{}.},label=lst:lab:l19:run]
# 1. start PRG32-QT on this computer (it listens on port 8080)

# 2. in the PRG32 folder, with ESP-IDF active: build the cartridge
python3 -m prg32 cartridge build ~/arc12/code/l19.S --portable \
    --entry-prefix l19 --name l19 --out l19.prg32

# 3. send it to PRG32-QT and start it
python3 -m prg32 esp32c6 upload-and-run l19.prg32 \
    --url http://localhost:8080 --slot cart0

# 4. look, press, compare with the previous slide
\end{shell}
